\chapter{Meining utan målestokk}

\begin{motto}
Faget oppdaga at forma tyder noko.\\ Så gløymde det å spørje kven som avgjer kva ho tyder.
\end{motto}

Ein kunne tru at det å slå fast at form følgjer funksjon var tom, ville tvinge faget til å finne ein betre grunn å stå på. På sett og vis gjorde det òg det, og dette kapitlet handlar om forsøket. Frå syttitalet og gjennom åttitalet skjedde det ei vending i korleis faget tenkte om kva ei form i det heile var til for. Funksjonen var avslørt som eit slør; men i staden for å rive ned sløret og spørje kva som verkeleg verkar på forma, bytte faget berre ut sløret med eit nytt og finare. Den nye grunnen heitte \emph{meining}. Forma følgjer ikkje lenger funksjonen, sa den nye tenkinga; forma \emph{tyder} noko, ho kommuniserer, ho ber innhald, og oppgåva til formgjevaren er å gjere dette innhaldet rett. Vendinga var i utgangspunktet både sann og frigjerande. Problemet var at faget gjorde med meininga nett det det hadde gjort med funksjonen: det sette henne i sentrum utan å skaffe seg ein målestokk for henne, og dermed bytte det berre ein udefinert grunn mot ein annan.

\section{Frå funksjon til teikn}

For å forstå kor naturleg vendinga kom, må ein sjå kva som var i ferd med å rakne. Modernismen hadde lova at forma kunne lesast rett ut av funksjonen, reint, utan stil og utan retorikk. Den lovnaden var slått hol i, både teoretisk og estetisk; den kvite, sakssvarande, ornamentlause forma som hevda å vere fråvêret av smak, var sjølv blitt ein smak, og ein som folk var i ferd med å bli grundig leie. Då Robert Venturi skreiv at han føretrekte det samansette framfor det reine, det tvitydige framfor det eintydige, og at «mindre er ein kjedsemd» som svar på modernismens «mindre er meir», markerte han eit oppgjer som var i emning over heile faget \autocite{venturi1966}. Forma fekk lov til å vere fleirtydig att, til å spele på historie og kontekst, til å seie meir enn éin ting. Då han nokre år seinare, saman med Denise Scott Brown og Steven Izenour, gjekk til Las Vegas for å studere skilt og fasadar som meiningsberande system heller enn som vulgaritet å forakte, var grepet fullbyrda: bygningen var ikkje lenger fyrst og fremst ein lausing på eit funksjonelt problem, han var eit teikn, noko som kommuniserer \autocite{venturi1972learning}. Arkitekturen las Las Vegas, og lærte å lese seg sjølv som språk.

I produktdesign tok same rørsla form som det som vart kalla \emph{produktsemantikk}. Klaus Krippendorff og Reinhart Butter sette omgrepet i omløp tidleg på åttitalet med ein enkel og slåande påstand: ein gjenstand sin viktigaste eigenskap er ikkje kva han gjer, men kva han \emph{tyder} for den som brukar han \autocite{krippendorff1984semantics}. Ein kontrollpanel kommuniserer kva slags maskin det høyrer til; ein telefon seier noko om kven du er; ei form ber symbolske kvalitetar som styrer korleis folk forstår og brukar henne. Designarens oppgåve var difor ikkje berre å løyse eit teknisk problem, men å forme tydinga. Og rundt om i verda braut dette laus som ein heil estetikk. I Milano sette Ettore Sottsass og Memphis-gruppa farge, mønster og sitat på ein måte som var medvite meiningsfull og medvite ufunksjonell; ein bokhylle skulle ikkje fyrst og fremst halde bøker, ho skulle seie noko. Dette var ikkje slurv. Det var ein tese: at forma talar, og at funksjonen alltid hadde vore eit påskot for ikkje å høyre etter kva ho sa.

Ein bør sjå kor djupt denne rørsla var samanvoven med tida si tankesett, for ho var ikkje berre eit motefenomen i form. Ho fall saman med ei brei intellektuell vending der språket, teiknet og tolkinga vart sett i sentrum av nær sagt alle humanistiske fag; semiotikken lærte ein heil generasjon å sjå alt som tekst, alt som system av teikn som kunne lesast. At produktdesign og arkitektur tok opp same vokabularet — gjenstanden som teikn, forma som ytring, brukaren som lesar — var ikkje tilfeldig. Det var faget si deltaking i ei allmenn dreiing frå substans til betydning. Og for designfaget var dreiinga dobbelt freistande, fordi ho lova å gje faget noko å snakke om som verken var den utskjelte funksjonen eller den ærleg innrømte smaken. Å seie at ein form «tyder» noko, lydde langt meir alvorleg, langt meir intellektuelt, enn å seie at ein likte henne. Den semantiske vendinga gav faget eit vakabular som låt teoretisk. Spørsmålet boka må stille er om vokabularet bar ein teori, eller berre lydnaden av ein.

\section{Det vendinga fekk rett i}

Lat oss vere generøse, for vendinga fortener det, og generøsiteten gjer kritikken skarpare. Den semantiske vendinga hadde rett i noko viktig, noko funksjonalismen hadde fortrengt og som faget framleis har godt av å hugse. Gjenstandar \emph{tyder} verkeleg noko. Mennesket har aldri omgått tinga sine som rein nytte; vi les status, tilhøyrsle, identitet og minne inn i og ut av forma, og ein teori om design som ikkje har plass til dette, er ein teori om ein menneskeart som ikkje finst. Då Donald Norman seinare argumenterte for at den emosjonelle responsen på ein gjenstand — gleda, tilknytinga, det vakre — ikkje er eit overflatisk tillegg til funksjonen men ein eigen og legitim dimensjon ved korleis ting verkar på oss, gav han den semantiske intuisjonen ei psykologisk underbygging \autocite{norman2004emotional}. Vi elskar og hatar tinga våre, og det er ikkje irrasjonelt; det er ein del av kva tinga er for oss. Vendinga frå funksjon til meining var difor ei reell utviding av kva faget kunne sjå. Penny Sparke har vist korleis heile det tjuande hundreåret sin designkultur ikkje let seg forstå utan smak, identitet og symbolsk verdi som drivkrefter — at gjenstanden alltid var eit kulturelt teikn, ikkje berre eit verktøy \autocite{sparke1986}. Den semantiske vendinga gjorde berre eksplisitt det designhistoria heile tida hadde vist.

Så langt er alt vel. Faget tok meininga på alvor, og det var rett å gjere. Men hugs kva eit fundament skal gjere. Eit fundament skal ikkje berre peike på den rette dimensjonen; det skal gje deg ein måte å avgjere ting innanfor den dimensjonen. Funksjonalismen sin feil var ikkje at funksjon var feil dimensjon — funksjon er ein heilt verkeleg ting som verkar på forma — men at formelen ikkje kunne brukast til å avgjere noko, fordi «funksjon» kunne tyde alt. Spørsmålet vi må stille den semantiske vendinga er difor presis det same spørsmålet: gav ho faget ein måte å avgjere?

\section{Kven avgjer kva ei form tyder?}

Her sviktar det, og det sviktar på ein måte som er strukturelt identisk med det vi alt har sett. Sei at forma tyder noko. Straks reiser to spørsmål seg, og faget svarar på korkje det eine eller det andre. Det fyrste: \emph{kven avgjer} kva ei form tyder? Designaren har ein intensjon, men tydinga oppstår ikkje i designaren; ho oppstår i møtet mellom gjenstanden og den som ser. Brukaren les si eiga tyding, forma av sin eigen kultur, si eiga tid, si eiga erfaring, og denne tydinga treng ikkje ha noko å gjere med det designaren la inn. Sottsass kunne meine at hylla seier ein ting; ho som ser henne, kan lese ei heilt anna, eller ingenting. Meininga er ikkje ein eigenskap forma har, slik vekt er; ho er ein relasjon som vert oppretta i lesinga, og difor utanfor designarens kontroll. Det andre spørsmålet er endå hardare: \emph{korleis veit ein at ein har rett?} Når funksjonen var grunnen, kunne ein i det minste prøve om stolen bar. Når meininga er grunnen, kva er prøva? Korleis avgjer ein om ei form lukkast i å tyde det ho skal? Det finst inga måling, ingen terskel, ingen instans som kan seie at lesinga er rett eller gal. Det finst berre tolkingar, og blant tolkingar finst det ingen domstol.

Ein kunne svare at dette er for strengt, at all kommunikasjon lever med tolkingsrom og likevel fungerer; vi forstår kvarandre i grove trekk trass i at orda er fleirtydige. Det er sant, men sjå kvifor det ikkje reddar faget. Språket fungerer fordi det har konvensjonar som er etablerte, delte og lærte — ein ordbok, ein grammatikk, eit fellesskap som korrigerer den som brukar orda feil. Designfaget har ingen ordbok over former. Det finst ingen etablert, delt kode som seier kva ei viss kurve, ein viss materiale, ein viss proporsjon tyder, og difor ingen instans som kan seie at ein designar las koden feil. Der språket har konvensjon og korrigering, har faget berre den einskilde designaren sin påstand om at forma tyder slik, sett opp mot den einskilde brukaren sin lesnad. Det er ikkje kommunikasjon i den forstand som fungerer; det er to private tolkingar som tilfeldigvis møtest ved same gjenstand. Faget tok lånet frå semiotikken — gjenstanden som teikn — men lét vere å skaffe det semiotikken byggjer på: ein delt kode. Det fekk teiknet og dropp koden, og eit teikn utan kode tyder ingenting bestemt for nokon.

Klaus Krippendorff, som meir enn nokon prøvde å gjere produktsemantikken til ei sann grunngjeving for faget, såg dette og freista å svare. I sitt store seinverk reformulerte han heile prosjektet og kalla det utan blygsel \emph{ein ny grunnvoll for design}: design er framføre alt det å gje ting meining, og faget burde organiserast kring meiningsskapinga heller enn kring forma eller funksjonen \autocite{krippendorff2006}. Innsatsen er sjølve det boka leitar etter — eit forsøk på å reise eit fundament — og difor er det dømmande at også her vert grunnen verande utan målestokk. Krippendorff flytter meininga frå designarens hovud til brukarens og til språket og praksisen kring gjenstanden, kva han kallar meininga sin avhengnad av interessentane; men dette gjer berre problemet meir akutt, for no finst tydinga overalt og difor ingen stad i særleg forstand. Når meininga ligg hjå alle som omgår gjenstanden, kan han bety alt for nokon, og då har faget framleis ingen måte å seie at ei tyding er meir rett enn ei anna, eller at ein designar har gjort arbeidet sitt betre enn ein annan. Den nye grunnvollen er like udefinert som den gamle. Han har berre flytta udefinertheita frå funksjon til meining, og kalla flyttinga eit framsteg.

Sjå då strukturen reint, for han er den same som i kapitlet om den tomme formelen. «Forma følgjer funksjonen» var tom fordi funksjon kunne tyde alt. «Forma ber meining» er tom på akkurat same vis, fordi meining kan vere kva som helst og avgjerast av kven som helst. I begge tilfella har faget identifisert ein verkeleg dimensjon ved forma — fyrst nytte, så tyding — og forveksla det å peike på dimensjonen med det å ha ein teori om henne. Ein teori om meining i forma måtte ha sagt kva som gjer ei tyding rettare enn ei anna, korleis tydingar oppstår av formtrekk og kontekst på måtar ein kan føreseie og prøve, og kven sin lesnad som tel og kvifor. Ingenting av dette kom. Det som kom, var ein lisens til å kalle kva som helst designaren føretrekte for «meiningsfullt» — den same lisensen funksjonalismen gav til å kalle kva som helst «funksjonelt». Den semantiske vendinga avskaffa ikkje den tomme formelen. Ho gjorde han elegant.

Det er denne elegansen som gjer vendinga så lett å forveksle med framsteg, og difor verdt å avsløre eit augneblink til. Funksjonalismens tomleik var grov og lett å gjennomskode når nokon fyrst peika på henne; «funksjon» som tyder alt, er ein ganske naken sirkel. Meiningas tomleik er finare, fordi ho kjem kledd i eit heilt teoriapparat — teikn, kode, tolking, interessent — som ber preg av lærdom og difor avleier mistanken. Eit fag som seier at det studerer korleis gjenstandar produserer meining for ulike interessentar i ulike kontekstar, lyder som om det driv forsking. Men under det lærde laget ligg den same tomme strukturen: ingen målestokk, ingen domstol, ingen måte å ta feil på. Og hugs frå tidlegare i boka kva det tyder at eit fag ikkje kan ta feil — det tyder at det heller ikkje kan lære. Den semantiske vendinga gav faget eit rikare vokabular til å skildre kva forma gjer, og ikkje eit einaste nytt verktøy til å avgjere om ho gjer det godt. Han som no skulle felle dommen «denne forma tyder rett», stod like hjelpelaus som han som før sa «denne forma er funksjonell» — berre meir veltalande hjelpelaus.

\section{Når meining vert til vare}

Det er ein grunn til at dette ikkje berre er ein logisk skavank, men ein farleg ein, og det er her kapitlet peikar framover. Ein dimensjon utan målestokk er ein dimensjon utan motstand. Når det ikkje finst nokon måte å seie at ei tyding er feil, finst det heller ingen måte å seie at ei tyding er \emph{misvisande} — og misvisande tyding har eit anna namn når nokon tener på henne: manipulasjon. Deyan Sudjic har sagt dette så reint at det er vanskeleg å misforstå: tinga me omgjev oss med er eit språk, og det språket vert i aukande grad skrive ikkje av oss, men \emph{for} oss, av dei som vil at me skal ynskje, kjøpe, høyre til \autocite{sudjic2008}. Forma talar; men når forma talar utan ein målestokk for sanning, kan ho lyge like lett som ho kan seie sant, og faget har ingen instans å skilje dei med. Den semantiske vendinga gav faget makt til å skape meining og fråtok det samstundes evna til å seie noko om kva meininga burde tene. Ho rusta opp forma som kommunikasjon og let spørsmålet om kommunikasjonen var ærleg eller utnyttande vere uråd å stille innanfrå faget.

Dette er ikkje eit lite poeng å avslutte på; det er kimen til den andre halvdelen av boka. For så lenge meininga budde i statussymbol og Memphis-hyller, var innsatsen avgrensa: det stod om smak og pengar, om luksus og forfengelegheit, ting eit fag kan leve med å vere uklart om. Men idet forma flytter inn på skjermen, og tydinga ho ber vert til signal som kan målast, testast og optimaliserast mot åtferd, vert ein udefinert meiningsgrunn til noko langt meir potent. Då vert meining til ein reiskap for å styre handling, og formgjevaren til den som byggjer styringa. Vendinga frå funksjon til meining ser uskuldig ut så lenge ho gjeld ein stol. Ho ser annleis ut når objektet ho gjev meining er ein brukar.

\vignett

For meininga skulle snart få ein heilt ny stad å bu, og ein heilt ny målestokk — ikkje sanning, men effekt. Forma vart til grensesnitt, brukaren vart til måltavle, og faget fann endeleg det som såg ut som ein objektiv grunn å stå på. Det er der vi går no: til skjermen, og til mannen som lærte faget å sjå brukaren.

\vidarelesing{\fullcite{krippendorff2006}; \fullcite{krippendorff1984semantics}; \fullcite{venturi1966}; \fullcite{venturi1972learning}; \fullcite{sudjic2008}; \fullcite{sparke1986}; \fullcite{norman2004emotional}.}
